%% ──────────────────────────────────────────────
%%  فصل ۱۳ – فعالیت‌های میسیونری و ایدئولوژیک
%%  فایل: chapters/ch13-missionary-ideological.tex
%% ──────────────────────────────────────────────
\chapter{فعالیت‌های میسیونری و ایدئولوژیک: تسهیل‌گری فرهنگی مداخله}
\label{ch:missionary}

%% ── چکیده‌ی فصل ──
\begin{tcolorbox}[mybox, title={چکیده‌ی فصل}]
فصل‌های پیشین ابعاد تاریخی (۴–۱۰)، روان‌شناختی (۱۱)، و اقتصادی-جامعه‌شناختی (۱۲) نجات‌طلبی و مداخله را بررسی کردند. اما بُعد سومی نیز وجود دارد که بدون آن تحلیل ناقص خواهد ماند: \concept{بُعد فرهنگی-ایدئولوژیک}. مداخله‌ی خارجی تقریباً همیشه با یک \concept{پروژه‌ی فرهنگی} همراه بوده است — پروژه‌ای که ذهن‌ها را پیش از سرزمین‌ها فتح می‌کند. از \concept{میسیونرهای مسیحی} در آفریقا و آسیا (قرن‌های ۱۶ تا ۲۰) گرفته تا \concept{صدور انقلاب اسلامی ایران}، از \concept{پان‌اسلاویسم} روسی و \concept{پان‌ترکیسم} عثمانی تا \concept{پان‌عربیسم} ناصری، از \concept{انترناسیونالیسم کمونیستی} شوروی تا \concept{دموکراسی‌سازی} آمریکایی — همه‌ی اینها اَشکال مختلف یک پدیده‌ی واحدند: \concept{تسهیل‌گری فرهنگی مداخله}. این فصل هر یک از این پروژه‌ها را بررسی، مقایسه، و در چارچوب مدل شش‌وجهی تحلیل می‌کند.
\end{tcolorbox}

%% ================================================================
\section{درآمد: قدرت نرم و فتح ذهن‌ها}
\label{sec:ch13-intro}
%% ================================================================

\needspace{6\baselineskip}
\thinker{جوزف نای} (\lat{Joseph Nye})، نظریه‌پرداز روابط بین‌الملل، در سال ۱۹۹۰ مفهوم \concept{قدرت نرم} (\lat{Soft Power}) را مطرح کرد: توانایی شکل‌دهی به ترجیحات دیگران از طریق جذابیت فرهنگی، ارزش‌های سیاسی، و سیاست خارجی — نه زور یا تطمیع.\footnote{\lr{Nye, Joseph S. Jr. \textit{Soft Power: The Means to Success in World Politics}. New York: PublicAffairs, 2004.}} اما آنچه نای «قدرت نرم» می‌نامد، تاریخی بسیار طولانی‌تر از اصطلاحش دارد: از زمانی که نخستین امپراتوری‌ها تشکیل شدند، فتح فرهنگی پیش‌درآمد، همراه، یا مکمل فتح نظامی بوده است.

\thinker{آنتونیو گرامشی} (\lat{Antonio Gramsci})، فیلسوف مارکسیست ایتالیایی، مفهوم \concept{هژمونی فرهنگی} (\lat{Cultural Hegemony}) را برای توضیح این پدیده به‌کار برد: طبقه‌ی حاکم نه فقط از طریق زور، بلکه از طریق شکل‌دهی به فرهنگ، ارزش‌ها، و «عقل سلیم» جامعه سلطه‌ی خود را حفظ می‌کند — و این سلطه به‌قدری «طبیعی» جلوه می‌کند که فرودستان آن را درونی می‌سازند.\footnote{\lr{Gramsci, Antonio. \textit{Selections from the Prison Notebooks}. Edited and translated by Quintin Hoare and Geoffrey Nowell Smith. New York: International Publishers, 1971, pp. 12–13, 57–58.}}

\begin{tcolorbox}[defbox, title={مفهوم کلیدی: تسهیل‌گری فرهنگی مداخله}]
\concept{تسهیل‌گری فرهنگی مداخله} (\lat{Cultural Facilitation of Intervention}) به مجموعه‌ی فعالیت‌های فرهنگی، مذهبی، ایدئولوژیک، آموزشی، و رسانه‌ای‌ای اطلاق می‌شود که — خواسته یا ناخواسته — زمینه‌ی ذهنی و اجتماعی را برای پذیرش، دعوت، یا مشروعیت‌بخشی به مداخله‌ی خارجی فراهم می‌آورند. این فعالیت‌ها ممکن است:
\begin{itemize}
\item \textbf{پیشگام} مداخله باشند: میسیونرها پیش از استعمارگران می‌رسند
\item \textbf{همزمان} با مداخله باشند: تبلیغات جنگی
\item \textbf{پیامد} مداخله باشند: «ملت‌سازی فرهنگی» پس از اشغال
\item یا \textbf{مستقل} از مداخله‌ی نظامی باشند: صدور ایدئولوژی بدون حمله‌ی نظامی
\end{itemize}
\end{tcolorbox}

همان‌گونه که در فصل~\ref{ch:social-psychology} دیدیم، سوگیری‌های شناختی و هویت اجتماعی زمینه‌ی روان‌شناختی نجات‌طلبی را فراهم می‌کنند؛ و در فصل~\ref{ch:sociology-economics} دیدیم که ساختارهای اقتصادی-سیاسی زمینه‌ی مادی آن را ایجاد می‌نمایند. اکنون می‌بینیم که پروژه‌های فرهنگی-ایدئولوژیک، بُعد سوم — \concept{زمینه‌ی معنایی} — را تکمیل می‌کنند: به مداخله «معنا» می‌بخشند و آن را از «تجاوز» به «رسالت» تبدیل می‌نمایند.


%% ================================================================
\section{میسیونرهای مسیحی: پیشقراولان استعمار}
\label{sec:ch13-christian-missionaries}
%% ================================================================

\needspace{6\baselineskip}
\subsection{تاریخچه: از حواریون تا استعمار}
\label{subsec:ch13-missionaries-history}

فعالیت میسیونری مسیحی سابقه‌ای دوهزارساله دارد، اما شکل سازمان‌یافته و مرتبط با قدرت سیاسی آن عمدتاً از قرن شانزدهم — همزمان با آغاز عصر اکتشافات و استعمار اروپایی — شکل گرفت. فرمان پاپی \concept{اینتر سیتیرا} (\lat{Inter Caetera}، ۱۴۹۳) از پاپ \histref{الکساندر ششم}{پاپ} به اسپانیا و پرتغال حق «تبشیر و تمدن‌سازی» در سرزمین‌های «کشف‌شده» را اعطا کرد — و بدین ترتیب تبشیر مسیحی با استعمار پیوند رسمی خورد.\footnote{\lr{Mignolo, Walter D. \textit{The Darker Side of the Renaissance: Literacy, Territoriality, and Colonization}. 2nd ed., Ann Arbor: University of Michigan Press, 2003, pp. 29–67.}}

\thinker{ادوارد سعید} (\lat{Edward Said}) در کتاب تأثیرگذار \textit{شرق‌شناسی} نشان داده است که پروژه‌ی دانش/قدرت استعماری — شامل مطالعات شرق‌شناسانه، نقشه‌برداری، مردم‌شناسی، و تبلیغ مذهبی — ابزاری برای «شناختن برای تسلط‌یافتن» بود.\footnote{\lr{Said, Edward W. \textit{Orientalism}. New York: Pantheon Books, 1978.}}

\subsection{میسیونرها در آفریقا}
\label{subsec:ch13-africa-missionaries}

در آفریقا، رابطه‌ی میسیونرها و استعمار پیچیده و چندلایه بود. \thinker{نگوگی وا تیونگو} (\lat{Ngũgĩ wa Thiong'o})، نویسنده‌ی کنیایی، در کتاب \textit{استعمارزدایی از ذهن} استدلال کرده است که بزرگ‌ترین دستاورد استعمار نه فتح سرزمین‌ها، بلکه \concept{فتح ذهن‌ها} بود — و میسیونرها ابزار اصلی این فتح بودند:\footnote{\lr{wa Thiong'o, Ngũgĩ. \textit{Decolonising the Mind: The Politics of Language in African Literature}. London: James Currey, 1986.}}

\begin{tcolorbox}[quotebox, title={نقل‌قول تاریخی}]
«گلوله‌ی تفنگ، قلب استعمارزده را هدف قرار می‌داد. اما گلوله‌ی فرهنگی، ذهنش را هدف قرار می‌داد. و فتح ذهن، فتح نهایی بود.»\\
— \thinker{نگوگی وا تیونگو} (\lat{Ngũgĩ wa Thiong'o})، \textit{استعمارزدایی از ذهن}، ۱۹۸۶\footnote{\lr{wa Thiong'o, Ngũgĩ. \textit{Decolonising the Mind}, op. cit., p. 9.}}
\end{tcolorbox}

میسیونرها در آفریقا سه کارکرد اصلی داشتند:

\begin{enumerate}
\item \textbf{کارکرد فرهنگی:} تبدیل زبان، مذهب، و ارزش‌های بومی. آموزش به زبان استعمارگر، تخریب نظام‌های دانش بومی، معرفی مفاهیم «تمدن» و «وحشیگری».\footnote{\lr{Comaroff, Jean and John L. Comaroff. \textit{Of Revelation and Revolution, Vol. 1: Christianity, Colonialism, and Consciousness in South Africa}. Chicago: University of Chicago Press, 1991.}}

\item \textbf{کارکرد اطلاعاتی:} میسیونرها اطلاعات ارزشمندی درباره‌ی جغرافیا، جمعیت، ساختارهای قبیله‌ای، و زبان‌های محلی جمع‌آوری می‌کردند که مستقیماً در خدمت اداره‌ی استعماری قرار می‌گرفت.

\item \textbf{کارکرد مشروعیت‌بخشی:} حضور میسیونرها به استعمار «چهره‌ی انسانی» می‌بخشید: استعمارگران فقط برای غارت نیامده‌اند — برای «نجات روح‌ها» و «تمدن‌سازی» نیز آمده‌اند.
\end{enumerate}

\begin{tcolorbox}[casebox, title={مطالعه‌ی موردی: کنگوی بلژیک — میسیونرها و استعمار}]
کنگوی بلژیک (۱۹۰۸–۱۹۶۰) نمونه‌ی بارز همکاری میسیونرها و استعمار بود. \histref{لئوپولد دوم}{پادشاه بلژیک} کنگو را به‌عنوان «مالکیت شخصی» اداره می‌کرد و سیستم کار اجباری برای استخراج لاستیک و عاج ایجاد کرد که میلیون‌ها نفر را کشت. میسیونرهای کاتولیک — هرچند برخی از آن‌ها از خشونت انتقاد کردند — در مجموع بخشی از ساختار استعماری بودند: مدارس میسیونری تنها نظام آموزشی بودند و فارغ‌التحصیلانشان «نخبگان دست‌آموز» (\concept{évolués}) استعمار شدند.\footnote{\lr{Hochschild, Adam. \textit{King Leopold's Ghost: A Story of Greed, Terror, and Heroism in Colonial Africa}. Boston: Houghton Mifflin, 1998, pp. 128–149.}}

اما — و این نکته مهم است — همین نظام آموزشی میسیونری، رهبران استقلال‌طلبی چون \histref{پاتریس لومومبا}{نخست‌وزیر کنگو} را نیز پرورش داد. تناقض میسیونرها: ابزار سلطه‌ای که گاهی ابزار رهایی نیز می‌شد.
\end{tcolorbox}

\subsection{میسیونرها در ایران و خاورمیانه}
\label{subsec:ch13-iran-missionaries}

در ایران، فعالیت میسیونری مسیحی — به‌ویژه پروتستان‌های آمریکایی — از اوایل قرن نوزدهم آغاز شد. میسیونرهای آمریکایی مدارس، بیمارستان‌ها، و مطابع تأسیس کردند (از جمله کالج آمریکایی تهران و ارومیه). اگرچه تعداد تغییر مذهب‌دادگان ناچیز بود، تأثیر فرهنگی آن‌ها — به‌ویژه در معرفی آموزش مدرن، پزشکی غربی، و چاپ — قابل‌توجه بود.\footnote{\lr{Zirinsky, Michael P. "Render Therefore unto Caesar the Things Which Are Caesar's: American Presbyterian Educators and Reza Shah," \textit{Iranian Studies}, Vol. 26, No. 3-4, 1993, pp. 337–356.}}

\thinker{جان الدر} (\lat{John Elder})، یکی از آخرین میسیونرهای آمریکایی در ایران، در خاطراتش اعتراف کرده است که میسیونرها — هرچند خوش‌نیت — ناخواسته به ایجاد تصویر «آمریکای خیّر» در ذهن ایرانیان کمک کردند؛ تصویری که بعدها کودتای ۲۸ مرداد ۱۳۳۲ آن را نابود کرد.\footnote{\lr{Elder, John. \textit{History of the American Presbyterian Mission in Iran, 1834–1960}. Tehran: Literature Committee of the Church Council of Iran, 1960.}}


%% ================================================================
\section{صدور انقلاب اسلامی ایران}
\label{sec:ch13-iran-revolution-export}
%% ================================================================

\needspace{6\baselineskip}
انقلاب اسلامی ۱۹۷۹ ایران نه‌فقط یک تحول داخلی، بلکه پروژه‌ای ایدئولوژیک با ابعاد بین‌المللی بود. \histref{آیت‌الله خمینی}{رهبر انقلاب اسلامی} بارها تأکید کرد که انقلاب «صادراتی» است و نباید به مرزهای ایران محدود شود:

\begin{tcolorbox}[quotebox, title={نقل‌قول تاریخی}]
«ما باید انقلاب‌مان را به تمام جهان صادر کنیم... این تنها یک انقلاب ایرانی نیست، بلکه انقلاب اسلامی است.»\\
— \histref{آیت‌الله خمینی}{رهبر انقلاب اسلامی ایران}\footnote{\lr{Khomeini, Ruhollah. "New Year's Message," 21 March 1980. In \textit{Sahifeh-ye Imam}, Vol. 12. Tehran: Institute for Compilation and Publication of Imam Khomeini's Works, pp. 56–58.}}
\end{tcolorbox}

\subsection{مکانیزم‌های صدور انقلاب}
\label{subsec:ch13-export-mechanisms}

صدور انقلاب اسلامی از چندین مسیر انجام شد:

\begin{tcolorbox}[wavebox, title={مسیرهای صدور انقلاب اسلامی}]
\begin{enumerate}
\item \textbf{نمونه‌سازی} (\lat{Exemplarism}): خود انقلاب به‌عنوان الگو — نشان دادن اینکه یک ملت مسلمان می‌تواند ابرقدرت وابسته‌اش را سرنگون کند

\item \textbf{حمایت از گروه‌های همسو:}
\begin{itemize}
\item تأسیس و تقویت حزب‌الله لبنان (از ۱۹۸۲)\footnote{\lr{Norton, Augustus Richard. \textit{Hezbollah: A Short History}. Princeton: Princeton University Press, 2007.}}
\item حمایت از گروه‌های شیعی عراقی (حزب الدعوه، مجلس اعلا)
\item حمایت از حوثی‌های یمن (هرچند حوثی‌ها زیدی‌اند، نه دوازده‌امامی)
\item ارتباط با حماس و جهاد اسلامی فلسطین (سنّی، اما مورد حمایت ایران)
\end{itemize}

\item \textbf{دیپلماسی فرهنگی:}
\begin{itemize}
\item سازمان فرهنگ و ارتباطات اسلامی
\item کنفرانس بین‌المللی وحدت اسلامی
\item بورسیه‌ی تحصیلی برای طلاب خارجی در قم و مشهد
\end{itemize}

\item \textbf{رسانه:} شبکه‌ی العالم (عربی)، پرس‌تی‌وی (انگلیسی)، هیسپان‌تی‌وی (اسپانیایی)

\item \textbf{محور مقاومت:} شکل‌دهی به ائتلاف ژئوپولیتیک ایران-سوریه-حزب‌الله-حماس-حوثی‌ها
\end{enumerate}
\end{tcolorbox}

\subsection{ارزیابی: موفقیت یا شکست صدور انقلاب؟}
\label{subsec:ch13-export-assessment}

\thinker{ولی نصر} (\lat{Vali Nasr}) در کتاب \textit{احیای شیعی} استدلال کرده است که صدور انقلاب در شکل اصلی‌اش (الهام‌بخشی به انقلاب‌های اسلامی در دیگر کشورها) عمدتاً شکست خورده، اما در شکل ژئوپولیتیک‌اش (ایجاد شبکه‌ی نفوذ منطقه‌ای) نسبتاً موفق بوده.\footnote{\lr{Nasr, Vali. \textit{The Shia Revival: How Conflicts Within Islam Will Shape the Future}. New York: W.W. Norton, 2006.}}

\begin{tcolorbox}[failbox, title={محدودیت‌های صدور انقلاب اسلامی}]
\begin{enumerate}
\item \textbf{مرز فرقه‌ای:} جذابیت اصلی انقلاب برای شیعیان بود — اکثریت مسلمانان جهان سنّی‌اند
\item \textbf{ناسیونالیسم:} حتی شیعیان عراقی در جنگ ایران-عراق علیه ایران جنگیدند
\item \textbf{مدل اقتصادی ناموفق:} جمهوری اسلامی الگوی اقتصادی جذابی ارائه نکرد
\item \textbf{سرکوب داخلی:} خبرهای سرکوب آزادی‌های مدنی جذابیت مدل را کاهش داد
\item \textbf{جنگ‌های فرقه‌ای:} حمایت ایران از اسد در سوریه — با هزینه‌ی جانی بالای سنّی‌ها — تصویر ایران را در جهان عرب خدشه‌دار کرد
\end{enumerate}
\end{tcolorbox}


%% ================================================================
\section{پان‌اسلاویسم: روسیه و رسالت اسلاوی}
\label{sec:ch13-panslavism}
%% ================================================================

\needspace{6\baselineskip}
\concept{پان‌اسلاویسم} (\lat{Pan-Slavism}) جنبشی ایدئولوژیک بود که وحدت همه‌ی ملت‌های اسلاو — تحت رهبری یا حمایت روسیه — را تبلیغ می‌کرد. این ایدئولوژی در قرن نوزدهم به‌ویژه در بالکان نقش مهمی در مشروعیت‌بخشی به مداخلات روسیه ایفا کرد.

\thinker{نیکلای دانیلوسکی} (\lat{Nikolay Danilevsky}) در کتاب \textit{روسیه و اروپا} (۱۸۶۹) چارچوب نظری پان‌اسلاویسم را ارائه کرد: تمدن اسلاوی تمدنی مجزا و برتر از تمدن غربی (ژرمنی-رومنی) است و روسیه رسالت رهبری آن را دارد.\footnote{\lr{Danilevsky, Nikolay. \textit{Russia and Europe: A Look at the Cultural and Political Relations of the Slavic World to the Romano-Germanic World}. 1869. English trans. Stephen Woodburn. Bloomington: Slavica Publishers, 2013.}}

\subsection{پان‌اسلاویسم و مسئله‌ی شرق}
\label{subsec:ch13-panslavism-eastern}

همان‌گونه که در فصل~\ref{ch:european} بحث شد، \concept{مسئله‌ی شرق} (\lat{Eastern Question}) — سرنوشت سرزمین‌های عثمانی در بالکان — عرصه‌ی اصلی رقابت قدرت‌های اروپایی در قرن نوزدهم بود. روسیه از پان‌اسلاویسم به‌عنوان ابزار مشروعیت‌بخشی به مداخلاتش استفاده کرد:

\begin{tcolorbox}[casebox, title={پان‌اسلاویسم به‌مثابه‌ی ابزار مداخله}]
\begin{itemize}
\item \textbf{جنگ روسیه-عثمانی ۱۸۷۷–۱۸۷۸:} روسیه به بهانه‌ی «حمایت از برادران اسلاو» و «آزادی بلغارها» به عثمانی حمله کرد. نتیجه: ایجاد بلغارستان مستقل (عملاً تحت نفوذ روسیه)\footnote{\lr{Todorova, Maria. \textit{Imagining the Balkans}. Oxford: Oxford University Press, 1997, pp. 76–88.}}
\item \textbf{حمایت از صربستان:} روسیه از اواخر قرن نوزدهم حامی اصلی صربستان بود — حمایتی که مستقیماً به جنگ جهانی اول (۱۹۱۴) منجر شد
\item \textbf{پان‌اسلاویسم مدرن:} پوتین اگرچه از اصطلاح «پان‌اسلاویسم» استفاده نمی‌کند، اما مفهوم \concept{«جهان روسی»} (\lat{Russkiy Mir}) — حمایت از روس‌زبانان و فرهنگ روسی در «خارج نزدیک» — بازخوانی مدرن همان ایدئولوژی است\footnote{\lr{Laruelle, Marlène. \textit{Russian Nationalism: Imaginaries, Doctrines, and Political Battlefields}. London: Routledge, 2019, pp. 148–172.}}
\end{itemize}
\end{tcolorbox}


%% ================================================================
\section{پان‌ترکیسم: از عثمانی تا آسیای مرکزی}
\label{sec:ch13-panturkism}
%% ================================================================

\needspace{6\baselineskip}
\concept{پان‌ترکیسم} (\lat{Pan-Turkism}) ایدئولوژی وحدت همه‌ی ملل ترک‌زبان — از ترکیه تا آسیای مرکزی و شینجیانگ چین — بود. این ایدئولوژی در اواخر دوره‌ی عثمانی توسط متفکرانی چون \thinker{ضیاء گوکالپ} (\lat{Ziya Gökalp}) و \thinker{اسماعیل گاسپرالی} (\lat{İsmail Gasprinsky}) توسعه یافت.\footnote{\lr{Landau, Jacob M. \textit{Pan-Turkism: From Irredentism to Cooperation}. 2nd ed., Bloomington: Indiana University Press, 1995.}}

\subsection{پان‌ترکیسم و ایران}
\label{subsec:ch13-panturkism-iran}

پان‌ترکیسم برای ایران اهمیت ویژه‌ای دارد — به‌ویژه در ارتباط با مسئله‌ی آذربایجان. همان‌گونه که در فصل~\ref{ch:modern-iran} بحث شد، جمهوری خودمختار آذربایجان (۱۹۴۵–۱۹۴۶) تحت حمایت شوروی، تجربه‌ای از تقاطع پان‌ترکیسم و ژئوپولیتیک بود.

در دوره‌ی معاصر، \histref{رجب طیب اردوغان}{رئیس‌جمهور ترکیه} سازمان \concept{کشورهای ترک‌زبان} (\lat{Organization of Turkic States}) را تقویت کرده و رویکرد \concept{نوعثمانی‌گرایی} (\lat{Neo-Ottomanism}) را دنبال می‌کند که عناصری از پان‌ترکیسم را با اسلام‌گرایی ترکیبی و ناسیونالیسم ترکی درمی‌آمیزد.\footnote{\lr{Taşpınar, Ömer. "Turkey's Middle East Policies: Between Neo-Ottomanism and Kemalism," \textit{Carnegie Endowment Papers}, No. 10, 2008.}}

\begin{tcolorbox}[enemybox, title={هشدار: پان‌ترکیسم و تمامیت ارضی ایران}]
پان‌ترکیسم یکی از تهدیدات بالقوه برای تمامیت ارضی ایران محسوب می‌شود. فعالیت رسانه‌ها و شبکه‌های اجتماعی ترک‌زبان، تأسیس کانال‌های تلویزیونی مرتبط (گوناز تی‌وی)، و حمایت برخی محافل ترکیه و جمهوری آذربایجان از جریان‌های قومی تجزیه‌طلب، نمونه‌هایی از \concept{تسهیل‌گری فرهنگی} بالقوه مداخله‌اند.\footnote{\lr{Shaffer, Brenda. \textit{Borders and Brethren: Iran and the Challenge of Azerbaijani Identity}. Cambridge, MA: MIT Press, 2002.}}

با این حال، باید از اغراق نیز پرهیز کرد: اکثریت آذربایجانی‌های ایران هویت ملی ایرانی دارند و جریان‌های تجزیه‌طلب حمایت مردمی گسترده‌ای ندارند. نکته‌ی مهم آن است که بحران‌های داخلی (اقتصادی، سیاسی) می‌تواند فضا را برای بهره‌برداری خارجی از شکاف‌های قومی باز کند — مکانیزمی که در فصل~\ref{ch:social-psychology} (بازتعریف مرز خودی/بیگانه) بحث شد.
\end{tcolorbox}


%% ================================================================
\section{پان‌عربیسم و ناصریسم}
\label{sec:ch13-panarabism}
%% ================================================================

\needspace{6\baselineskip}
\concept{پان‌عربیسم} (\lat{Pan-Arabism}) ایدئولوژی وحدت همه‌ی عرب‌ها بر اساس زبان، فرهنگ، و تاریخ مشترک بود. این ایدئولوژی در قرن بیستم — به‌ویژه با ظهور \histref{جمال عبدالناصر}{رئیس‌جمهور مصر، ۱۹۵۴–۱۹۷۰} — به نیروی سیاسی قدرتمندی تبدیل شد.

\thinker{باسل دیویدسون} (\lat{Bassam Tibi}) در تحلیل خود از پان‌عربیسم، آن را نمونه‌ای از \concept{ناسیونالیسم پان‌ها} (\lat{Pan-Nationalisms}) می‌داند: ایدئولوژی‌هایی که مرزهای موجود دولتی را نامشروع تلقی می‌کنند و وحدت بر اساس هویت فراملی (قومی، زبانی، مذهبی) را تبلیغ می‌نمایند.\footnote{\lr{Tibi, Bassam. \textit{Arab Nationalism: Between Islam and the Nation-State}. 3rd ed., London: Macmillan, 1997.}}

\subsection{ناصریسم: وعده و شکست}
\label{subsec:ch13-nasserism}

ناصر با ملی‌سازی کانال سوئز (۱۹۵۶) و مقاومت در برابر تجاوز سه‌گانه‌ی بریتانیا-فرانسه-اسرائیل، به قهرمان جهان عرب تبدیل شد. اما اتحاد مصر و سوریه (\concept{جمهوری متحده‌ی عربی}، ۱۹۵۸–۱۹۶۱) شکست خورد و شکست در جنگ شش‌روزه‌ی ۱۹۶۷ پان‌عربیسم را عملاً نابود کرد.\footnote{\lr{Ajami, Fouad. \textit{The Arab Predicament: Arab Political Thought and Practice since 1967}. Updated ed., Cambridge: Cambridge University Press, 1992.}}

\begin{tcolorbox}[wavebox, title={پان‌عربیسم و مداخله}]
پان‌عربیسم به‌عنوان ابزار مداخله عمل کرد:
\begin{itemize}
\item \textbf{مداخله‌ی مصر در یمن (۱۹۶۲–۱۹۷۰):} ناصر به بهانه‌ی حمایت از «انقلاب جمهوری‌خواهانه» در یمن شمالی مداخله‌ی نظامی کرد — «ویتنام مصر» که ۲۶٬۰۰۰ سرباز مصری تلفات داشت\footnote{\lr{Ferris, Jesse. \textit{Nasser's Gamble: How Intervention in Yemen Caused the Six-Day War and the Decline of Egyptian Power}. Princeton: Princeton University Press, 2013.}}
\item \textbf{کودتاهای عربی:} ایدئولوژی پان‌عربیسم توجیه‌گر بسیاری از کودتاها بود: سوریه (۱۹۶۳)، عراق (۱۹۶۸)، لیبی (۱۹۶۹)
\item \textbf{صدام و کویت:} صدام حمله به کویت (۱۹۹۰) را با زبان پان‌عربیستی توجیه کرد: «بازگرداندن شاخه‌ی جداشده» به تنه‌ی اصلی
\end{itemize}
\end{tcolorbox}

پان‌عربیسم اکنون عمدتاً مرده است، اما ایده‌ی «وحدت عربی» همچنان ابزار رتوریکی مفیدی برای توجیه مداخلات (مثلاً مداخله‌ی عربستان در یمن) باقی مانده است.


%% ================================================================
\section{انترناسیونالیسم کمونیستی: شوروی و رسالت جهانی}
\label{sec:ch13-communism}
%% ================================================================

\needspace{6\baselineskip}
\concept{انترناسیونالیسم کمونیستی} (\lat{Communist Internationalism}) شاید بزرگ‌ترین و سازمان‌یافته‌ترین پروژه‌ی صدور ایدئولوژی در قرن بیستم بود. از تأسیس \concept{کمینترن} (\lat{Comintern}، ۱۹۱۹) تا فروپاشی شوروی (۱۹۹۱)، اتحاد جماهیر شوروی به‌طور سیستماتیک تلاش کرد انقلاب کمونیستی را به سراسر جهان صادر کند.\footnote{\lr{Westad, Odd Arne. \textit{The Global Cold War: Third World Interventions and the Making of Our Times}. Cambridge: Cambridge University Press, 2005.}}

\subsection{مکانیزم‌های صدور}
\label{subsec:ch13-communism-mechanisms}

\begin{enumerate}
\item \textbf{احزاب کمونیست محلی:} شوروی در تقریباً هر کشوری حزب کمونیست ایجاد یا حمایت کرد — از حزب توده‌ی ایران تا حزب کمونیست اندونزی
\item \textbf{آموزش کادرها:} دانشگاه \concept{پاتریس لومومبا} در مسکو (تأسیس ۱۹۶۰) هزاران «انقلابی» از جهان سوم را آموزش داد\footnote{\lr{Katsakioris, Constantin. "Creating a Socialist Intelligentsia: Soviet Educational Aid and Its Impact on Africa (1960–1991)," \textit{Cahiers d'études africaines}, No. 226, 2017, pp. 259–288.}}
\item \textbf{کمک نظامی:} تسلیحات، مشاوران نظامی، آموزش چریکی — از کوبا و ویتنام تا آنگولا و اتیوپی
\item \textbf{رسانه و تبلیغات:} رادیو مسکو به ده‌ها زبان پخش می‌شد
\item \textbf{مداخله‌ی نظامی مستقیم:} مجارستان (۱۹۵۶)، چکسلواکی (۱۹۶۸)، افغانستان (۱۹۷۹)
\end{enumerate}

\subsection{دکترین برژنف: «حاکمیت محدود»}
\label{subsec:ch13-brezhnev}

\concept{دکترین برژنف} (\lat{Brezhnev Doctrine}) — که پس از تهاجم به چکسلواکی (۱۹۶۸) صورت‌بندی شد — اعلام کرد که «اردوگاه سوسیالیستی» حق دارد در هر کشور عضو که «دستاوردهای سوسیالیسم» تهدید شود مداخله کند. این دکترین عملاً مفهوم \concept{«حاکمیت محدود»} (\lat{Limited Sovereignty}) را مطرح کرد — و شباهت آشکاری با مداخلات آمریکایی به بهانه‌ی «حفاظت از دموکراسی» داشت.\footnote{\lr{Ouimet, Matthew J. \textit{The Rise and Fall of the Brezhnev Doctrine in Soviet Foreign Policy}. Chapel Hill: University of North Carolina Press, 2003.}}

\begin{tcolorbox}[wavebox, title={تقارن‌های ایدئولوژیک: شوروی و آمریکا}]
\begin{itemize}
\item \textbf{شوروی:} «جهان باید سوسیالیستی شود» → مداخله برای «حفاظت از دستاوردهای سوسیالیسم»
\item \textbf{آمریکا:} «جهان باید دموکراتیک شود» → مداخله برای «حفاظت از دموکراسی و آزادی»
\item \textbf{هر دو:} ادعای جهانشمولی ارزش‌ها + حق مداخله برای حفاظت از آن ارزش‌ها + سرکوب مخالفان داخلی «خودی‌ها» + حمایت از دیکتاتورهای «خودی» در جهان سوم
\end{itemize}
تفاوت اصلی: شوروی صراحتاً مداخله را حق خود می‌دانست (دکترین برژنف)؛ آمریکا مداخله را «انتخاب آزاد ملت‌ها» جلوه می‌داد (اما در عمل مشابه عمل می‌کرد).
\end{tcolorbox}


%% ================================================================
\section{«دموکراسی‌سازی» آمریکایی: NED، USAID، و قدرت نرم}
\label{sec:ch13-democracy-promotion}
%% ================================================================

\needspace{6\baselineskip}
\subsection{ساختار نهادی ترویج دموکراسی}
\label{subsec:ch13-ned-usaid}

پس از فروپاشی شوروی، ایالات متحده پروژه‌ی \concept{ترویج دموکراسی} (\lat{Democracy Promotion}) را با شدت بیشتری دنبال کرد. ساختار نهادی این پروژه شامل:

\begin{tcolorbox}[defbox, title={نهادهای اصلی ترویج دموکراسی آمریکایی}]
\begin{itemize}
\item \concept{بنیاد ملی دموکراسی} (\lat{National Endowment for Democracy — NED}): تأسیس ۱۹۸۳ به فرمان ریگان. بودجه‌ی سالانه: حدود ۳۰۰ میلیون دلار (۲۰۲۲). بنیاد NED آشکارا آنچه را سیا قبلاً پنهانی انجام می‌داد، به‌صورت علنی انجام می‌دهد.\footnote{\lr{Allen Weinstein, co-founder of NED, quoted in David Ignatius, "Innocence Abroad: The New World of Spyless Coups," \textit{The Washington Post}, 22 September 1991: "A lot of what we do today was done covertly 25 years ago by the CIA."}}
\item \concept{آژانس توسعه‌ی بین‌المللی آمریکا} (\lat{USAID}): بودجه‌ی سالانه بیش از ۲۵ میلیارد دلار — شامل برنامه‌های «حکمرانی خوب» و «جامعه‌ی مدنی»
\item \concept{انستیتوی جمهوری‌خواه بین‌المللی} (\lat{IRI}) و \concept{انستیتوی دموکراتیک ملی} (\lat{NDI}): وابسته به دو حزب اصلی آمریکا
\item \concept{خانه‌ی آزادی} (\lat{Freedom House}): رتبه‌بندی «آزادی» کشورها — مورد استفاده‌ی سیاست‌گذاران
\item رسانه‌ها: رادیو فردا/آزادی (\lat{RFE/RL})، صدای آمریکا (\lat{VOA})، رادیو سوا (\lat{Radio Sawa})
\end{itemize}
\end{tcolorbox}

\subsection{نقد ترویج دموکراسی}
\label{subsec:ch13-democracy-critique}

\thinker{ویلیام بلوم} (\lat{William Blum}) در کتاب \textit{دولت سرکش} مستند کرده است که آمریکا بین ۱۹۴۵ و ۲۰۰۵ در بیش از ۵۰ کشور تلاش کرده دولت‌ها را سرنگون کند، در بیش از ۳۰ کشور در انتخابات دخالت نموده، و در بیش از ۲۰ کشور ترور سیاسی انجام داده — همه به نام «دموکراسی» و «آزادی».\footnote{\lr{Blum, William. \textit{Rogue State: A Guide to the World's Only Superpower}. 3rd ed., Monroe, ME: Common Courage Press, 2005.}}

\thinker{نیکلاس گیلهات} (\lat{Nicolas Guilhot}) در تحلیل جامعی از NED نشان داده است که «ترویج دموکراسی» اغلب به معنای ترویج نوع خاصی از دموکراسی (لیبرال، بازارمحور، غربی) است — نه هر دموکراسی‌ای. دموکراسی‌های «نامطلوب» (مثلاً حماس در غزه ۲۰۰۶ یا مرسی در مصر ۲۰۱۲) به رسمیت شناخته نمی‌شوند.\footnote{\lr{Guilhot, Nicolas. \textit{The Democracy Makers: Human Rights and the Politics of Global Order}. New York: Columbia University Press, 2005.}}

\begin{tcolorbox}[casebox, title={مطالعه‌ی موردی: «انقلاب‌های رنگی»}]
\concept{انقلاب‌های رنگی} (\lat{Color Revolutions}) — گرجستان (انقلاب رز ۲۰۰۳)، اوکراین (انقلاب نارنجی ۲۰۰۴)، قرقیزستان (انقلاب لاله ۲۰۰۵) — نمونه‌هایی از تقاطع جنبش‌های مردمی واقعی با حمایت سازمان‌یافته‌ی غربی بودند.

\textbf{نقش NED و جورج سوروس:} بنیاد \concept{جامعه‌ی باز} (\lat{Open Society Foundations}) جورج سوروس و NED از سازمان‌های مدنی، رسانه‌های مستقل، و ناظران انتخاباتی در این کشورها حمایت مالی و فنی کردند.\footnote{\lr{Mitchell, Lincoln A. \textit{The Color Revolutions}. Philadelphia: University of Pennsylvania Press, 2012.}}

\textbf{ارزیابی:}
\begin{itemize}
\item این جنبش‌ها بدون نارضایتی واقعی مردمی ممکن نبودند — نقش مردم واقعی بود
\item اما حمایت خارجی (مالی، آموزشی، رسانه‌ای) نقش تسهیل‌گری مهمی داشت
\item نتایج متفاوت: گرجستان → دموکراسی نسبی؛ اوکراین → بی‌ثباتی و سپس تهاجم روسیه؛ قرقیزستان → بازگشت اقتدارگرایی
\item روسیه این انقلاب‌ها را «کودتاهای غربی» تلقی کرد — و واکنشش مستقیماً به بحران اوکراین ۲۰۱۴ و ۲۰۲۲ منجر شد
\end{itemize}
\end{tcolorbox}


%% ================================================================
\section{جدول تطبیقی: پروژه‌های ایدئولوژیک مداخله}
\label{sec:ch13-comparative-table}
%% ================================================================

\needspace{6\baselineskip}

\begin{landscape}
\begin{table}[htbp]
\centering
\caption{جدول تطبیقی پروژه‌های ایدئولوژیک-فرهنگی تسهیل‌گر مداخله}
\label{tab:ch13-comparative}
\begin{adjustbox}{max width=\linewidth}
\begin{tabularx}{\linewidth}{
    >{\raggedleft\arraybackslash}p{2.2cm}
    >{\raggedleft\arraybackslash}p{2cm}
    >{\centering\arraybackslash}p{1.8cm}
    >{\raggedleft\arraybackslash}p{2.5cm}
    >{\raggedleft\arraybackslash}p{2.5cm}
    >{\raggedleft\arraybackslash}p{2.5cm}
    >{\centering\arraybackslash}p{1.5cm}
}
\toprule
\rowcolor{PrimaryDeep}
\textcolor{white}{\textbf{پروژه}} &
\textcolor{white}{\textbf{عامل اصلی}} &
\textcolor{white}{\textbf{دوره}} &
\textcolor{white}{\textbf{ابزار اصلی}} &
\textcolor{white}{\textbf{جغرافیای هدف}} &
\textcolor{white}{\textbf{رابطه با مداخله‌ی نظامی}} &
\textcolor{white}{\textbf{موفقیت}} \\
\midrule
تبشیر مسیحی & کلیساها + دولت‌های استعماری & قرن ۱۶–۲۰ & مدرسه، بیمارستان، کتاب مقدس & آفریقا، آسیا، اقیانوسیه & پیشگام و همراه استعمار & بالا \\
\rowcolor{NeutralLight}
صدور انقلاب اسلامی & ایران & ۱۹۷۹–کنون & حمایت از گروه‌ها، رسانه، بورسیه & خاورمیانه، شیعیان جهان & همراه (سوریه) و مستقل & متوسط \\
پان‌اسلاویسم & روسیه & قرن ۱۹–۲۱ & حمایت نظامی، فرهنگی، مذهبی & بالکان، اوکراین & مشروعیت‌بخش مداخله‌ی نظامی & بالا \\
\rowcolor{NeutralLight}
پان‌ترکیسم & عثمانی/ترکیه & قرن ۱۹–۲۱ & رسانه، آموزش، دیپلماسی فرهنگی & قفقاز، آسیای مرکزی، ایران & محدود (آذربایجان ۲۰۲۰) & متوسط \\
پان‌عربیسم & مصر ناصری & ۱۹۵۰–۱۹۷۰ & رادیو، کودتا، اتحاد سیاسی & جهان عرب & مستقیم (یمن ۱۹۶۲) & شکست \\
\rowcolor{NeutralLight}
انترناسیونالیسم کمونیستی & شوروی/چین & ۱۹۱۹–۱۹۹۱ & احزاب محلی، آموزش، سلاح & جهانی & مستقیم (مجارستان، افغانستان) & متوسط→شکست \\
دموکراسی‌سازی & آمریکا & ۱۹۸۳–کنون & NED، USAID، رسانه، NGO & جهانی & هم‌راستا (عراق، لیبی) & متفاوت \\
\rowcolor{NeutralLight}
وهابیسم/سلفی‌گری & عربستان & ۱۹۷۰–کنون & مساجد، مدارس، کتاب، رسانه & جهان اسلام & غیرمستقیم (جهادی‌ها) & بالا \\
\bottomrule
\end{tabularx}
\end{adjustbox}
\end{table}
\end{landscape}


%% ================================================================
\section{وهابیسم و سلفی‌گری صادراتی عربستان}
\label{sec:ch13-wahhabism}
%% ================================================================

\needspace{6\baselineskip}
یکی از مهم‌ترین و کمتر بحث‌شده‌ی پروژه‌های صدور ایدئولوژی، گسترش جهانی \concept{وهابیسم/سلفی‌گری} توسط عربستان سعودی است. از دهه‌ی ۱۹۷۰ — پس از افزایش قیمت نفت — عربستان سرمایه‌گذاری عظیمی در ساخت مساجد، مدارس دینی (\concept{مدرسه}ها)، چاپ و توزیع رایگان کتب مذهبی، و تربیت روحانیون سلفی در سراسر جهان اسلام (و فراتر از آن) انجام داده است.\footnote{\lr{Lacroix, Stéphane. \textit{Awakening Islam: The Politics of Religious Dissent in Contemporary Saudi Arabia}. Cambridge, MA: Harvard University Press, 2011.}}

\thinker{ناصر الحزیمی} (\lat{Nasser al-Huzaimi}) و دیگر پژوهشگران سعودی تخمین زده‌اند که عربستان بین ۱۹۷۵ و ۲۰۰۵ بیش از ۱۰۰ میلیارد دلار صرف تبلیغ وهابیسم در جهان کرده است — بزرگ‌ترین پروژه‌ی تبلیغ ایدئولوژیک تاریخ.\footnote{\lr{Commins, David. \textit{The Wahhabi Mission and Saudi Arabia}. London: I.B. Tauris, 2006.}}

\begin{tcolorbox}[enemybox, title={پیامدهای صدور وهابیسم}]
\begin{itemize}
\item \textbf{رادیکالیزاسیون:} بسیاری از رهبران و اعضای القاعده، داعش، و گروه‌های جهادی در مدارس تأمین‌مالی‌شده‌ی سعودی آموزش دیدند
\item \textbf{فرقه‌گرایی:} تبلیغ وهابیسم تنش‌های فرقه‌ای (سنی-شیعه) را در سراسر جهان اسلام تشدید کرد
\item \

```latex
textbf{تخریب میراث فرهنگی:} ایدئولوژی وهابی بسیاری از بقاع متبرکه، مقبره‌ها، و آثار تاریخی اسلامی (حتی در مکه و مدینه) را تخریب کرده\footnote{\lr{Irfan al-Alawi and Stephen Schwartz. "Destruction of Islamic Heritage in Saudi Arabia," \textit{The Guardian}, 23 October 2012.}}
\item \textbf{همگن‌سازی:} تلاش برای جایگزینی اسلام‌های محلی و متکثر (صوفی، حنفی، شافعی) با قرائتی واحد و سخت‌گیرانه
\item \textbf{تغذیه‌ی تروریسم:} گزارش ۲۸ صفحه‌ی محرمانه‌ی گزارش کنگره درباره‌ی ۱۱ سپتامبر (منتشرشده در ۲۰۱۶) به ارتباطات مالی عربستان و عاملان حملات اشاره دارد\footnote{\lr{US Congress, Joint Inquiry into Intelligence Community Activities Before and After the Terrorist Attacks of September 11, 2001, "28 Pages," declassified July 2016.}}
\end{itemize}
\end{tcolorbox}

\thinker{کارن آرمسترانگ} (\lat{Karen Armstrong}) در تحلیل تاریخی خود تأکید کرده است که وهابیسم — که در قرن هجدهم یک جریان حاشیه‌ای و محلی در نجد بود — بدون ثروت نفتی و حمایت سیاسی دولت سعودی هرگز به ایدئولوژی جهانی تبدیل نمی‌شد. این نمونه‌ی روشنی از تقاطع \concept{نفرین منابع} (فصل~\ref{ch:sociology-economics}) با \concept{صدور ایدئولوژی} (این فصل) است.\footnote{\lr{Armstrong, Karen. "Wahhabism to ISIS: How Saudi Arabia Exported the Main Source of Global Terrorism," \textit{New Statesman}, 27 November 2014.}}


%% ================================================================
\section{نمودار ون: تقاطع پروژه‌های ایدئولوژیک}
\label{sec:ch13-venn}
%% ================================================================

\needspace{14\baselineskip}

\begin{figure}[htbp]
\centering
\begin{tikzpicture}[
    scale=0.85, transform shape
]

% دایره‌ها
\begin{scope}[opacity=0.3]
\fill[red!60] (0, 0) circle (3.5cm);
\fill[blue!60] (3, 0) circle (3.5cm);
\fill[green!50] (1.5, 2.6) circle (3.5cm);
\end{scope}

% حاشیه‌ی دایره‌ها
\draw[thick, red!70!black] (0, 0) circle (3.5cm);
\draw[thick, blue!70!black] (3, 0) circle (3.5cm);
\draw[thick, green!60!black] (1.5, 2.6) circle (3.5cm);

% برچسب دایره‌ها
\node[font=\small\bfseries, red!70!black] at (-1.8, -1.5) {\rl{مذهبی}};
\node[font=\tiny, red!70!black] at (-1.8, -2) {\rl{میسیونرها، وهابیسم،}};
\node[font=\tiny, red!70!black] at (-1.8, -2.4) {\rl{صدور انقلاب}};

\node[font=\small\bfseries, blue!70!black] at (4.8, -1.5) {\rl{ایدئولوژیک-سیاسی}};
\node[font=\tiny, blue!70!black] at (4.8, -2) {\rl{کمونیسم، دموکراسی‌سازی،}};
\node[font=\tiny, blue!70!black] at (4.8, -2.4) {\rl{پان‌عربیسم}};

\node[font=\small\bfseries, green!60!black] at (1.5, 5.3) {\rl{قومی-ناسیونالیستی}};
\node[font=\tiny, green!60!black] at (1.5, 4.8) {\rl{پان‌اسلاویسم، پان‌ترکیسم}};

% تقاطع‌ها
\node[font=\tiny, align=center] at (1.5, 0) {\rl{اسلام سیاسی}\\[2pt]\rl{(ایران، القاعده)}};
\node[font=\tiny, align=center] at (0.3, 1.5) {\rl{ارتدکسی}\\[2pt]\rl{+ اسلاویسم}\\[2pt]\rl{(روسیه)}};
\node[font=\tiny, align=center] at (2.7, 1.5) {\rl{ناسیونالیسم}\\[2pt]\rl{سکولار}\\[2pt]\rl{(ناصریسم)}};

% مرکز
\node[font=\tiny\bfseries, align=center, fill=yellow!20, rounded corners, inner sep=3pt] at (1.5, 0.9) {\rl{تسهیل‌گری}\\[2pt]\rl{فرهنگی مداخله}};

% عنوان
\node[font=\normalsize\bfseries, PrimaryDeep] at (1.5, -4) {\rl{تقاطع پروژه‌های ایدئولوژیک تسهیل‌گر مداخله}};

\end{tikzpicture}
\caption{نمودار ون: سه‌گانه‌ی پروژه‌های مذهبی، ایدئولوژیک-سیاسی، و قومی-ناسیونالیستی}
\label{fig:ch13-venn}
\end{figure}


%% ================================================================
\section{تحلیل تطبیقی: وجوه مشترک و متمایز}
\label{sec:ch13-comparative-analysis}
%% ================================================================

\needspace{6\baselineskip}
\subsection{وجوه مشترک}
\label{subsec:ch13-common}

بررسی تطبیقی هشت پروژه‌ی ایدئولوژیک-فرهنگی بررسی‌شده در این فصل، چند الگوی مشترک را آشکار می‌سازد:

\begin{tcolorbox}[wavebox, title={الگوهای مشترک پروژه‌های ایدئولوژیک}]
\textbf{۱. ادعای جهانشمولی:}
هر پروژه ادعا می‌کند ارزش‌هایش «جهانشمول»اند — نه خاص یک فرهنگ یا ملت: مسیحیت «حقیقت جهانی» است؛ کمونیسم «قانون تاریخ» است؛ دموکراسی لیبرال «پایان تاریخ» است؛ اسلام «دین فطرت» است. این ادعای جهانشمولی، مداخله را نه «تجاوز» بلکه «رسالت» جلوه می‌دهد.

\textbf{۲. شکاف میان شعار و عمل:}
در همه‌ی موارد، فاصله‌ی عمیقی میان ارزش‌های اعلام‌شده و عملکرد واقعی وجود دارد: میسیونرها عشق و صلح تبلیغ می‌کردند اما با استعمارگران همکاری داشتند؛ شوروی «رهایی ملل» را تبلیغ می‌کرد اما مجارستان و چکسلواکی را اشغال نمود؛ آمریکا «دموکراسی» صادر می‌کند اما از دیکتاتوری‌های متحد حمایت می‌نماید.

\textbf{۳. آماده‌سازی ذهنی برای مداخله:}
در اکثر موارد، فعالیت فرهنگی-ایدئولوژیک \textbf{پیش} از مداخله‌ی نظامی یا همزمان با آن صورت می‌گیرد: میسیونرها قبل از ارتش استعماری می‌رسند؛ رادیو مسکو قبل از تانک‌های شوروی آغاز به کار می‌کند؛ NED قبل از بمباران‌های ناتو فعال می‌شود.

\textbf{۴. ایجاد «نخبگان دست‌آموز»:}
هر پروژه‌ی ایدئولوژیک تلاش می‌کند گروهی از نخبگان محلی را پرورش دهد که ارزش‌های مداخله‌گر را درونی کرده و واسطه‌ی مشروعیت‌بخشی به مداخله شوند: \concept{évolués} استعمار، کادرهای حزب کمونیست، «فعالان جامعه‌ی مدنی» تربیت‌شده‌ی NED، طلاب وهابی تربیت‌شده‌ی عربستان.

\textbf{۵. بهره‌برداری از شکاف‌های داخلی:}
هر پروژه‌ی ایدئولوژیک از شکاف‌های موجود در جامعه‌ی هدف بهره‌برداری می‌کند: شکاف‌های قومی (پان‌ترکیسم)، فرقه‌ای (وهابیسم، صدور انقلاب)، طبقاتی (کمونیسم)، یا سیاسی (دموکراسی‌سازی).
\end{tcolorbox}

\subsection{وجوه متمایز}
\label{subsec:ch13-differences}

\begin{tcolorbox}[wavebox, title={تفاوت‌های کلیدی}]
\textbf{۱. جهت حرکت:}
\begin{itemize}
\item پروژه‌های مذهبی (میسیونرها، وهابیسم) عمدتاً از «بالا به پایین» عمل می‌کنند: از طریق نهادهای مذهبی
\item پروژه‌های سیاسی (کمونیسم) ادعای «پایین به بالا» دارند: از طریق بسیج طبقه‌ی کارگر
\item دموکراسی‌سازی ترکیبی است: نهادهای خارجی + فعالان محلی
\end{itemize}

\textbf{۲. رابطه با خشونت:}
\begin{itemize}
\item میسیونرهای مسیحی عمدتاً غیرخشونت‌آمیز (اما همراه با ارتش استعماری)
\item کمونیسم صراحتاً انقلاب خشونت‌آمیز را مشروع می‌دانست
\item وهابیسم/سلفی‌گری جهادی، خشونت را ابزار مقدس تلقی می‌کند
\item دموکراسی‌سازی ظاهراً غیرخشونت‌آمیز — اما اغلب با مداخله‌ی نظامی همراه می‌شود
\end{itemize}

\textbf{۳. بُرد زمانی:}
\begin{itemize}
\item تبشیر مسیحی: قرن‌ها — طولانی‌ترین پروژه
\item کمونیسم: هفت دهه (۱۹۱۷–۱۹۹۱) — کوتاه اما شدید
\item صدور انقلاب اسلامی: چهار دهه — ادامه‌دار
\item دموکراسی‌سازی آمریکایی: چهار دهه — ادامه‌دار اما تضعیف‌شده
\end{itemize}
\end{tcolorbox}


%% ================================================================
\section{مدل تحلیلی: سه لایه‌ی تسهیل‌گری فرهنگی مداخله}
\label{sec:ch13-model}
%% ================================================================

\needspace{14\baselineskip}

\begin{figure}[htbp]
\centering
\begin{tikzpicture}[
    >=Stealth,
    scale=0.82, transform shape,
    layer/.style={trapezium, draw=PrimaryDeep, fill=#1, trapezium angle=75,
        minimum width=5cm, minimum height=1.5cm, align=center, font=\small,
        text=black, trapezium stretches body},
]

% سه لایه (هرم)
\node[layer=red!20, minimum width=10cm] (top) at (0, 5) {\textbf{\rl{لایه‌ی ۱: تولید معنا}}\\[3pt]\rl{ایدئولوژی، مذهب، ارزش‌های «جهانشمول»}};

\node[layer=orange!20, minimum width=8cm] (mid) at (0, 3) {\textbf{\rl{لایه‌ی ۲: نهادسازی}}\\[3pt]\rl{مدرسه، رسانه، NGO، حزب، مسجد/کلیسا}};

\node[layer=yellow!20, minimum width=6cm] (bot) at (0, 1) {\textbf{\rl{لایه‌ی ۳: تربیت نخبگان}}\\[3pt]\rl{«évolués»، کادرها، فعالان، طلاب}};

% نتیجه
\node[rectangle, draw=PrimaryDeep, fill=PrimaryDeep!15, rounded corners,
    text width=5cm, minimum height=1.2cm, align=center, font=\small] (result)
    at (0, -1) {\textbf{\rl{نتیجه:}}\\[3pt]\rl{آماده‌سازی ذهنی برای}\\[2pt]\rl{پذیرش مداخله‌ی خارجی}};

% فلش‌ها
\draw[->, very thick, PrimaryDeep] (top.south) -- (mid.north);
\draw[->, very thick, PrimaryDeep] (mid.south) -- (bot.north);
\draw[->, very thick, red!70!black] (bot.south) -- (result.north);

% حاشیه‌نویسی
\node[font=\tiny, text=gray, right] at (5.5, 5) {\rl{چه باید باور کرد؟}};
\node[font=\tiny, text=gray, right] at (4.5, 3) {\rl{از چه مسیری باور منتقل شود؟}};
\node[font=\tiny, text=gray, right] at (3.5, 1) {\rl{چه کسانی باور را منتقل کنند؟}};

% عنوان
\node[font=\normalsize\bfseries, PrimaryDeep] at (0, 7) {\rl{مدل سه‌لایه‌ای تسهیل‌گری فرهنگی مداخله}};

\end{tikzpicture}
\caption{هرم سه‌لایه‌ای تسهیل‌گری فرهنگی مداخله}
\label{fig:ch13-pyramid}
\end{figure}


%% ================================================================
\section{مقاومت فرهنگی: ضد-هژمونی و استقلال فرهنگی}
\label{sec:ch13-resistance}
%% ================================================================

\needspace{6\baselineskip}
همان‌گونه که در فصل~\ref{ch:social-psychology} از «عاملیت جمعی» به‌عنوان نقطه‌ی خروج از چرخه‌ی روان‌شناختی نجات‌طلبی سخن گفتیم، در حوزه‌ی فرهنگی نیز اَشکال مقاومت وجود دارد.

\subsection{مفهوم ضد-هژمونی گرامشی}
\label{subsec:ch13-counter-hegemony}

\thinker{گرامشی} خود مفهوم \concept{ضد-هژمونی} (\lat{Counter-Hegemony}) را مطرح کرد: تلاش گروه‌های فرودست برای ایجاد فرهنگ، ارزش‌ها، و نهادهای جایگزین که هژمونی فرهنگی طبقه‌ی حاکم را به چالش بکشد. در زمینه‌ی بین‌المللی، ضد-هژمونی به معنای تولید روایت‌ها، مدل‌ها، و نهادهای فرهنگی مستقل از مرکز است.\footnote{\lr{Carroll, William K. (ed.). \textit{Organizing the 1\%: How Corporate Power Works}. Halifax: Fernwood Publishing, 2017, pp. 183–201.}}

\subsection{نمونه‌های مقاومت فرهنگی}
\label{subsec:ch13-cultural-resistance}

\begin{tcolorbox}[successbox, title={نمونه‌های مقاومت فرهنگی در برابر تسهیل‌گری مداخله}]
\begin{enumerate}
\item \textbf{جنبش نگریتود} (\lat{Négritude}): جنبش ادبی-فرهنگی آفریقایی (دهه‌ی ۱۹۳۰–۱۹۶۰) به رهبری \thinker{امه سزر} (\lat{Aimé Césaire}) و \thinker{لئوپولد سدار سنگور} (\lat{Léopold Sédar Senghor}) که ارزش فرهنگ‌های آفریقایی را در برابر هژمونی فرهنگی استعمار بازتعریف کرد\footnote{\lr{Césaire, Aimé. \textit{Discourse on Colonialism}. Trans. Joan Pinkham. New York: Monthly Review Press, 1972 (orig. 1950).}}

\item \textbf{الهیات رهایی‌بخش} (\lat{Liberation Theology}): جنبش الهیاتی آمریکای لاتین (دهه‌ی ۱۹۶۰–۱۹۸۰) که مسیحیت را نه ابزار سلطه، بلکه ابزار رهایی فقرا تفسیر کرد — بازتصاحب دین از دست استعمارگران\footnote{\lr{Gutiérrez, Gustavo. \textit{A Theology of Liberation: History, Politics, and Salvation}. Trans. Sister Caridad Inda and John Eagleson. Maryknoll, NY: Orbis Books, 1973.}}

\item \textbf{مکتب پسااستعماری} (\lat{Postcolonial Studies}): آثار \thinker{فرانتس فانون} (\lat{Frantz Fanon})، \thinker{ادوارد سعید}، \thinker{گایاتری اسپیواک} (\lat{Gayatri Spivak})، و \thinker{هومی بابا} (\lat{Homi Bhabha}) که ساختارهای دانش/قدرت استعماری را به چالش کشیدند

\item \textbf{جنبش عدم تعهد} (\lat{Non-Aligned Movement}): تلاش کشورهای جهان سوم در جنگ سرد برای ایجاد مسیر سومی بین شرق و غرب — هویت مستقل فرهنگی-سیاسی\footnote{\lr{Prashad, Vijay. \textit{The Darker Nations: A People's History of the Third World}. New York: The New Press, 2007.}}

\item \textbf{شبکه‌های رسانه‌ای جایگزین:} الجزیره (از ۱۹۹۶)، تلسور (\lat{teleSUR}) آمریکای لاتین، RT روسیه — تلاش برای شکستن انحصار روایتی غرب (هرچند خود ابزار دولت‌های دیگرند)

\item \textbf{آموزش انتقادی} (\lat{Critical Pedagogy}): \thinker{پائولو فریره} (\lat{Paulo Freire}) در \textit{آموزش ستمدیدگان} مدلی آموزشی ارائه کرد که ستمدیده را نه «موضوع نجات» بلکه «عامل رهایی خود» می‌بیند\footnote{\lr{Freire, Paulo. \textit{Pedagogy of the Oppressed}. Trans. Myra Bergman Ramos. New York: Continuum, 1970.}}
\end{enumerate}
\end{tcolorbox}


%% ================================================================
\section{جمع‌بندی فصل}
\label{sec:ch13-conclusion}
%% ================================================================

\needspace{6\baselineskip}
این فصل نشان داد که مداخله‌ی خارجی تقریباً همیشه با یک پروژه‌ی فرهنگی-ایدئولوژیک همراه بوده است — پروژه‌ای که ذهن‌ها را پیش از سرزمین‌ها یا همراه با آن‌ها فتح می‌کند. از میسیونرهای مسیحی در آفریقا تا NED آمریکایی در اوروپای شرقی، از صدور انقلاب اسلامی ایران تا صدور وهابیسم سعودی، الگوی واحدی دیده می‌شود: \textbf{تولید معنا → نهادسازی → تربیت نخبگان → آماده‌سازی ذهنی برای پذیرش مداخله.}

\begin{tcolorbox}[wavebox, title={یافته‌های کلیدی فصل ۱۳}]
\begin{enumerate}
\item \textbf{سه‌لایه‌ی تسهیل‌گری:} پروژه‌های ایدئولوژیک از سه لایه تشکیل شده‌اند: تولید معنا (ایدئولوژی)، نهادسازی (مدرسه، رسانه، NGO)، و تربیت نخبگان (کادرها و واسطه‌ها).

\item \textbf{ادعای جهانشمولی:} هر پروژه‌ی ایدئولوژیک ادعا می‌کند ارزش‌هایش جهانشمول‌اند — و بنابراین مداخله برای ترویج آن‌ها «رسالت» است، نه «تجاوز».

\item \textbf{شکاف شعار و عمل:} در همه‌ی موارد، فاصله‌ی عمیقی میان ارزش‌های اعلام‌شده و عملکرد واقعی وجود دارد. این شکاف، اگرچه از ناحیه‌ی مداخله‌گران نادیده گرفته می‌شود، در حافظه‌ی جمعی ملت‌های هدف ثبت و بازتولید می‌شود.

\item \textbf{بهره‌برداری از شکاف‌ها:} پروژه‌های ایدئولوژیک از شکاف‌های درونی جوامع هدف (قومی، فرقه‌ای، طبقاتی) بهره‌برداری می‌کنند — مکانیزمی که در فصل~\ref{ch:social-psychology} (بازتعریف مرز خودی/بیگانه) تحلیل شد.

\item \textbf{تقارن عملکردی:} پروژه‌های ظاهراً متضاد (کمونیسم و سرمایه‌داری، شیعه‌گری و وهابیسم، مسیحیت و اسلام) از نظر \textit{ساختار عملکردی} شباهت‌های بنیادین دارند: همه ادعای جهانشمولی می‌کنند، همه نخبگان دست‌آموز تربیت می‌کنند، و همه از مداخله حمایت یا با آن همراهی می‌نمایند.

\item \textbf{مقاومت فرهنگی ممکن است:} نگریتود، الهیات رهایی‌بخش، پسااستعمار، و آموزش انتقادی نشان می‌دهند که ملت‌ها می‌توانند در برابر هژمونی فرهنگی مقاومت کنند — به شرطی که روایت‌ها، نهادها، و نخبگان بومی مستقل داشته باشند.

\item \textbf{تکمیل مدل سه‌بعدی:} با ترکیب یافته‌های فصل‌های ۱۱ (روان‌شناسی)، ۱۲ (اقتصاد-جامعه‌شناسی)، و ۱۳ (فرهنگ-ایدئولوژی)، مدل جامعی از نجات‌طلبی شکل می‌گیرد: \textbf{ساختارهای اقتصادی} زمینه‌ی مادی، \textbf{مکانیزم‌های روان‌شناختی} زمینه‌ی ذهنی، و \textbf{پروژه‌های فرهنگی} زمینه‌ی معنایی مداخله را فراهم می‌آورند. فصل ۱۴ (جمع‌بندی) این سه بُعد را در یک مدل واحد ترکیب خواهد کرد.
\end{enumerate}
\end{tcolorbox}

\begin{tcolorbox}[quotebox, title={سخن پایانی فصل}]
«وقتی مبلّغان آمدند، ما زمین داشتیم و آن‌ها کتاب مقدس. به ما گفتند چشم‌ها را ببندیم و دعا کنیم. وقتی چشم‌ها را باز کردیم، ما کتاب مقدس داشتیم و آن‌ها زمین‌مان را.»\\
— ضرب‌المثل آفریقایی، منسوب به \thinker{جومو کنیاتا} (\lat{Jomo Kenyatta})\footnote{\lr{Kenyatta, Jomo. \textit{Facing Mount Kenya}. London: Secker and Warburg, 1938; often paraphrased as quoted above. See Sundkler, Bengt and Christopher Steed. \textit{A History of the Church in Africa}. Cambridge: Cambridge University Press, 2000, p. 73.}}\\[6pt]
«ایدئولوژی‌ها لباس‌هایی هستند که قدرت می‌پوشد تا برهنگی‌اش دیده نشود.»\\
— نویسنده
\end{tcolorbox}

\cleardoublepage